% 개념 25 85-34(인쇄 85쪽) — 확률밀도함수가 될 수 있는지 묻는 그래프: (0,1) 에서 (2,-1) 로 내려가는 선분과 점선 보조선.
% 스캔 화질이 낮아 TikZ 로 다시 그림. 축 0.7pt · 그래프 0.6pt · 라벨은 원문에 있는 것만(O · x · y · 1 · 2 · -1 · y=f(x)).
\begin{tikzpicture}[scale=0.75, line width=0.6pt]
  \draw[->, line width=0.7pt] (-0.45,0) -- (2.95,0) node[below, font=\small] {$x$};
  \draw[->, line width=0.7pt] (0,-1.75) -- (0,1.6) node[right, font=\small] {$y$};
  \node[below left=-2pt and -2pt, font=\small] at (0,0) {$\mathrm{O}$};
  \draw[dotted, line width=0.6pt] (2,0) -- (2,-1) -- (0,-1);
  \draw (0,1) -- (2,-1);
  \node[left, font=\small] at (0,1) {$1$};
  \node[left, font=\small] at (0,-1) {$-1$};
  \node[below, font=\small] at (1,0) {$1$};
  \node[above, font=\small] at (2,0) {$2$};
  \node[below, font=\small] at (1.75,-1.15) {$y=f(x)$};
\end{tikzpicture}
